%%%%%%%%%%%%%%%%%%%%%%%%%%%%%%%%%%%%%%%%%%%%%%%%%%%%%%%%%
% Para las listas
%%%%%%%%%%%%%%%%%%%%%%%%%%%%%%%%%%%%%%%%%%%%%%%%%%%%%%%%%

\newcommand{\cnumero}[2]{\tikz[baseline=(myanchor.base)]
\node[minimum size=0.2cm,circle,
inner sep=1pt,draw, #2,thick,fill=#2] (myanchor)
{\color{white}\bfseries\fontsize{8}{8}#1};}

\newcommand*{\itembolasazules}[1]{\protect\cnumero{#1}{azulF}}
\newcommand*{\itembolascolor}[2]{\protect\cnumero{#1}{#2}}
\newcommand*{\numberballgreen}{\protect\cnumero{\arabic*}{mygreen}}
\newcommand*{\numberballblue}{\protect\cnumero{\arabic*}{blue}}
\newcommand*{\numberballred}{\protect\cnumero{\arabic*}{red}}

\newcommand*{\numberballmDarkTeal}{\protect\cnumero{\arabic*}{mDarkTeal}}
\newcommand*{\numberballmDarkBrown}{\protect\cnumero{\arabic*}{mDarkBrown}}
\newcommand*{\numberballmLightBrown}{\protect\cnumero{\arabic*}{mLightBrown}}
\newcommand*{\numberballmLightGreen}{\protect\cnumero{\arabic*}{mLightGreen}}



%%%%%%%%%%%%%%%%%%%%%%%%%%%%%%%%%%%%%%%%%%%%%%%%%%%%%%%%%%%%%%%%%%%%%%%%%%%%%%%
% MIS Comandos Generales
%%%%%%%%%%%%%%%%%%%%%%%%%%%%%%%%%%%%%%%%%%%%%%%%%%%%%%%%%%%%%%%%%%%%%%%%%%%%%%%%
\newcommand{\blueremark}[1]{%
    \textcolor{blue}{\texttt{\textbf{#1}}}
    }
    
\newcommand{\rojo}[1]{\textcolor{red}{#1}}

\newcommand{\rojon}[1]{\textbf{\textcolor{red}{#1}}}

\newcommand{\azul}[1]{\textcolor{blue}{#1}}

\newcommand{\azuln}[1]{\textbf{\textcolor{blue}{#1}}}

\newcommand{\verde}[1]{\textcolor{green!50!black}{#1}}

\newcommand{\verden}[1]{\textbf{\textcolor{green!50!black}{#1}}}

\newcommand{\caja}[1]{\begin{center}%
    \begin{impt}%
        #1
    \end{impt}
    \end{center}}
    
\newcommand{\separ}{\\[2mm]}

\newcounter{ejercicio}
\setcounter{ejercicio}{0}

\newcounter{ejemplo}
\setcounter{ejemplo}{0}

\newcounter{ejerciciogeneral}
\setcounter{ejerciciogeneral}{0}

\newcommand{\ejerciciogeneral}[1]{\noindent\xmybox{\phantom{l}\textbf{Ejercicio \addtocounter{ejerciciogeneral}{1}\arabic{ejerciciogeneral} #1}\phantom{j}}}

\newcommand{\mybox}[1]{\noindent\xmybox{\phantom{l}\textbf{\addtocounter{ejercicio}{1}\arabic{ejercicio} #1}\phantom{j}}}

\newcommand{\marcador}[1]{\colorbox{yellow}{#1}}

\newcommand{\negrita}[1]{\textbf{#1}}
\newcommand{\subin}[1]{\textsubscript{#1}} 
\newcommand{\supin}[1]{\textsuperscript{#1}}
\newcommand{\negsubro}[1]{\negrita{\coloruline{red}{#1}}}

\newcommand{\letramayuscula}{\Alph*)}
\newcommand{\letraminuscula}{\textit{\textbf{\alph*})}}

%%%%%%%%%%%%%%%%%%%%%%%%%%%%%%%%%%%%%%%%%%%%%%%%%
% PARA SUBRAYAR REquiere paquete ulem
%%%%%%%%%%%%%%%%%%%%%%%%%%%%%%%%%%%%%%%%%%%%%%%%5
\makeatletter
    \newcommand{\coloruline}[2]{%
        \newcommand\temp@reduline{\bgroup\markoverwith
            {\textcolor{#1}{\rule[-0.5ex]{2pt}{0.4pt}}}\ULon}%
        \temp@reduline{#2}%
    }

    \newcommand{\coloruuline}[2]{%
        \UL@protected\def\temp@uuline{\leavevmode \bgroup
            \UL@setULdepth
            \ifx\UL@on\UL@onin \advance\ULdepth2.8\p@\fi
            \markoverwith{\textcolor{#1}{\lower\ULdepth\hbox
                {\kern-.03em\vbox{\hrule width.2em\kern1\p@\hrule}\kern-.03em}}}%
        \ULon}%
        \temp@uuline{#2}%
    }

    \newcommand{\coloruwave}[2]{%
        \UL@protected\def\temp@uwave{\leavevmode \bgroup
        \ifdim \ULdepth=\maxdimen \ULdepth 3.5\p@
        \else \advance\ULdepth2\p@
        \fi \markoverwith{\textcolor{#1}{\lower\ULdepth\hbox{\sixly \char58}}}\ULon}
        \font\sixly=lasy6 % does not re-load if already loaded, so no memory drain.
        \temp@uwave{#2}%
    }
\makeatother

%%%%%%%%%%%%%%%%%%%%%%%%%%%%%%%%%%%%%%%%%%%%%%%%%%%%%%%%%%%%%%%%%%%%%%%%%%%%%%%%

%%%%%%%%%%%%%%%%%%%%%%%%%%%%%%%%%%%%%%%%%%%%%%%%%%%%%%%%%%%%%%%%%%%%%%%%
% para hacer cancelto en rojo
%%%%%%%%%%%%%%%%%%%%%%%%%%%%%%%%%%%%%%%%%%%%%%%%%%%%%%%%%%%%%%%%%%%%%%%

\let\cancelorigcolor\CancelColor% Just for conveniency...

\newcommand{\CancelTo}[3][]{%
  \ifblank{#1}{}{%
    \renewcommand{\CancelColor}{#1}%
  }
  \cancelto{#2}{#3}% 
}


\newcommand{\canceltored}[2]{\CancelTo[\color{red}]{#1}{#2}}

\newcommand\cancelred[1]{\renewcommand\CancelColor{\color{red}}\cancel{#1}}


\newcounter{problema}
\setcounter{problema}{0}
\newcommand{\problema}[1]{\noindent\xmyboxorange{\phantom{l}\textbf{\addtocounter{problema}{1}\arabic{problema} #1}\phantom{j}}}

\newcommand{\ejercicio}[2]{\noindent\xmybox{\bfseries #1 Página #2}} %xmybox está en tcolorbox definida. Es para los ejercicos, indicando el ejercicio y la página.

%\newcommand{\mybox}[1]{\noindent\xmybox{\phantom{l}\textbf{#1}\phantom{j}}}

%\newcommand{\solution}[1]{\textcolor{red}{\textbf{\ul{Solución:} #1}}}

\newcommand{\solution}[1]{{\color{red}{\negrita{\ul{Solución:}} #1}}}

\newcommand{\datos}{\textbf{\ul{Datos:}}}


\newcommand{\ayuda}{\textbf{\ul{Ayuda:}}}




\newcounter{cuestion}
\setcounter{cuestion}{0}

\newcommand{\cuestion}[1]{\noindent\xmybox{\phantom{l}\textbf{Cuestión \addtocounter{cuestion}{1}\arabic{cuestion} #1}\phantom{j}}}


\newcommand{\cmmnt}[1]{\ignorespaces} % para comentar palabras en un texto.

\newcommand\myemptypage{
    \null
    \thispagestyle{empty}
    \addtocounter{page}{-1}
    \newpage
}


